 |m|\left(KR_N(\varepsilon)+\frac{250t_K^2}{\sqrt K}\right).    \tag{CB.18}
\] For bounded \(t_{K}\), \(KR_N=O(K^{1/4}/\sqrt N)\) follows directly from \(\varepsilon^2=t_K/(2s_N\sqrt K)\) and \(s_N\leq4\pi/N\). The remaining Gaussian characteristic-function error is \(o(1)\); no universal \(K^{-1/2}\) rate is asserted for that separate step.

\subsection{Microscopic consequences without conflating unfolding conventions}\label{coboundary:microscopic-consequences-without-conflating-unfolding-conventions}

For completeness we give the two spectral implications used in Theorem 2. First, the 4K Weyl atoms in \(C_{j}\) are distinct, so \(\operatorname{Tr}(\bar C^2)/N=4/K\). Distinctness follows within species from the unit conditions on \(1\pm h^j\) in the Pell construction. Between species, norm 1 versus -2 would require N to divide 3, which is impossible for \(K\geq3\). Consequently \[
 \tau_-(\bar C^2)\leq\frac{8N}{K(N-1)}.
\] The squared-overlap matrix between orthonormal eigenbases of two unitaries is doubly stochastic. Decomposing it into permutations yields an eigenvalue matching whose mean squared chord cost is bounded by normalized squared Hilbert--Schmidt distance. Circle distance in \(\mathbb R/\mathbb Z\) is at most one quarter of chord distance. Using this matching, (CB.7),(CB.14), and \([\widetilde U_0,\bar C]=0\), the RMS distance a of the spectrum of \(V_{\varepsilon}\) from the K-point lattice satisfies \[
 a\leq\frac{R_N(\varepsilon)}4+
          \frac{250\zeta^2}{\sqrt K}+
          \frac{|\zeta|}{2\pi}\sqrt{\frac{8N}{K(N-1)}}.         \tag{CB.19}
\] When \(t_{K}\) tends to zero, Ka tends to zero, and K/d tends to zero.

For an arbitrary d-point circle multiset with this RMS lattice distance, placing arcs of half-width \(w<1/(2K)\) around the lattice and using Markov's inequality gives, for its unfolded cyclic gaps \(S_{l}\), \[
 d^{-1}\#\{l:S_l>x\}\leq2a^2/w^2+K/d+2Kw/x.                  \tag{CB.20}
\] There are at most \(da^2/w^2\) bad points; at most twice as many gaps touch one, at most K other gaps leave an arc, and the within-arc gaps have total length at most 2Kw. This proves (CB.20). Partitioning the arcs into cells of length at most x/d and applying Cauchy--Schwarz to their occupancies also gives the ordered close-pair bound \[
 R_2(x)\geq\frac{\max(0,1-a^2/w^2)^2}{K/d+2Kw/x}-1.           \tag{CB.21}
\] Take \(w=(Ka)^{2/3}/K\) for (CB.20), and w=2a for (CB.21). These prove subcritical spacing collapse and divergent close-pair counts. The a=0 case follows by counting at most K nonzero gaps and using \(R_2(x)\geq d/K-1\). The average spacing remains exactly one: weak convergence to \(\delta_{0}\) does not imply convergence of first moments, which are carried by rare long gaps.

For bounded nonzero critical parameter a different argument is necessary. Let \(\phi_L(s)=L^{-1}(1-|s|/L)_+\), a compactly supported nonnegative function with integral one and Fourier transform \(\operatorname{sinc}^2(L\xi)\). For the eigenangles \(\theta_{l}\) in \(\mathbb R/\mathbb Z\), define \[
 C_d(\phi_L)=\frac1d\sum_{l\ne j}\sum_{q\in\mathbb Z}
                       \phi_L(d(\theta_l-\theta_j+q)).
\] Poisson summation, including the zero Fourier mode and the removed diagonal, proves, for \(d>L\), \[
 C_d(\phi_L)\geq1-\frac1L+
       2\sum_{m=1}^{M}\operatorname{sinc}^2(LmK/d)
                                 |\tau_-(V_\varepsilon^{mK})|^2. \tag{CB.22}
\] All omitted Fourier coefficients are nonnegative. Take K to infinity at fixed M, use (CB.2), then let M grow. For \(t>0\) this gives \[
 \liminf C_d(\phi_L)\geq\Theta(t)-1/L,\qquad
 \Theta(t)=\sum_{m\in\mathbb Z}e^{-m^2t^2}>1.                  \tag{CB.23}
\] For t=0 the lower bound is infinite. If the pair measures converged on compact tests to a bounded density R(s) tending to 1 at infinity, then \(\int\phi_LR\to1\) as L grows. This contradicts (CB.23). A bounded subsequence of \(|\varepsilon|K^{1/4}/\sqrt N\) has a further subsequence with finite limiting \(t_{K}\), so (CB.4) follows. This is a necessary-condition argument, not an exchange of a growing test support with the semiclassical limit.

An equivalent moving-window consequence is \[
 \liminf\operatorname{Var}_{\theta}\!\bigl[
      \#\{l:\theta_l\in[\theta,\theta+L/d)\}\bigr]
       \geq L^2(\Theta(t)-1).
\] Indeed the finite variance is \(L^2(C_d(\phi_L)-1)+L\). The underlying positive Fourier test was already established in the earlier Floquet note; it is not claimed as a new result of this pass.
